\documentclass[10pt,a4paper]{amsart}
\usepackage[margin=19mm]{geometry}
\usepackage{amsmath,amssymb,mathtools}
\usepackage[scr=rsfs,cal=euler]{mathalfa}
\usepackage{xfrac}
\usepackage[unicode,hidelinks,bookmarks=false]{hyperref}
\newcommand{\ie}{i.e.\ }
\newcommand{\cf}{cf.\ }
\newcommand{\resp}{respectively\ }
\newcommand{\Subsection}{\subsection}
\providecommand{\nobreakdash}{\nobreak\hbox{-}}
\setlength{\emergencystretch}{2em}
\multlinegap0pt
\usepackage{bm}
\usepackage{bbm}
\usepackage{dsfont}
\usepackage{xspace}
\usepackage{cancel}
\usepackage{mathtools}
\usepackage{amsthm}
\makeatletter
\@ifundefined{theorem}{\newtheorem{theorem}{Theorem}}{}
\@ifundefined{lemma}{\newtheorem{lemma}[theorem]{Lemma}}{}
\makeatother
\newcommand{\bbN}{{\ensuremath{\mathbb N}} }
\newcommand{\bbZ}{{\ensuremath{\mathbb Z}} }
\newcommand{\cC}{\ensuremath{\mathcal C}}
\newcommand{\pc}{\ensuremath{p_{\mathrm{c}} }}
\title[Catalan percolation]{Catalan percolation}
\author{Eleanor Archer, Ivailo Hartarsky, Brett Kolesnik, Sam Olesker-Taylor, Bruno Schapira, Daniel Valesin}
\date{Source: arXiv:2404.19583v2 (CC BY 4.0)}
\begin{document}
\maketitle

\noindent\emph{Source and licence.} Catalan percolation. Version arXiv:2404.19583v2 (CC BY 4.0), distributed under the Creative Commons Attribution 4.0 International licence, as recorded in the arXiv licence metadata for this exact version and shown on the abstract page https://arxiv.org/abs/2404.19583v2. Full rights notice: licenses/arxiv-2404.19583-CC-BY-4.0.txt.\par\smallskip

\noindent\textit{Editorial scope.} Counted unit: the Lemma whose source label is \texttt{lem:alpha:critical} in the article (source0.tex lines 1598--1601, source label \texttt{lem:alpha:critical}), delivered together with the article's own complete original proof of it (source0.tex lines 1602--1639, one \texttt{proof} environment of 38 lines). No mathematics is written, repaired or re-derived: statement and proof are the article's own text line for line. Required context retained uncounted: none. Every definition, display and label of those lines is retained unchanged. Omitted from this package and not redistributed: 1--1597, 1640--1810. Nothing inside a retained excerpt is omitted or altered: every retained body is delivered exactly as the article's lines are written, so no omission receipt of any kind accompanies this package. Citations: the retained excerpts carry no citation command at all, so no bibliography entry of the article is transcribed and no reference is redistributed. Reference adaptation: none was needed and none was made; every reference inside the delivered unit resolves inside it, so no unresolved reference or citation is produced. Printed slips kept literal: no printed word, symbol, hypothesis or number of the counted statement or of its proof was altered. Layout and class: the article's own class is \texttt{sn-jnl} with packages including graphicx, multirow, amsmath, amssymb, amsfonts, amsthm, mathrsfs, appendix, xcolor, textcomp, manyfoot, booktabs, algorithm, algorithmicx; the wrapper is the lane's pinned \texttt{amsart} editorial wrapper at 10pt on A4, which restates the theorem-like environments the retained text needs and adds the article's own macro definitions that the retained text needs (\texttt{\textbackslash bbN}, \texttt{\textbackslash bbZ}, \texttt{\textbackslash cC}, \texttt{\textbackslash pc}), with extra wrapper packages (bm, bbm, dsfont, xspace, cancel, mathtools). The delivered numbering is the unit's own. Assets: no retained span contains a figure environment, a graphics-inclusion command, a PDF-page inclusion, a table or a TikZ picture; no image file is copied into this package, the package declares no asset dependency and no image file is redistributed with it. Licence: version arXiv:2404.19583v2 of this article is distributed under the Creative Commons Attribution 4.0 International licence, as recorded in the arXiv licence metadata for this exact version and shown on the abstract page https://arxiv.org/abs/2404.19583v2. A full rights notice is delivered at licenses/arxiv-2404.19583-CC-BY-4.0.txt. Characters: every retained excerpt is pure ASCII, so the delivered engine, XeLaTeX with its default Latin Modern text font, reports no missing character.
\begin{lemma}
\label{lem:alpha:critical}
	If~$p > \pc(0)$, then~$\beta(p,0)^{-1} = \alpha(p,0) \ge 1$.
\end{lemma}

\begin{proof}
We write~$\cC_0 := \{(x,n){\in\bbZ^2}: x \in \xi_n(\{0\})\}$ {for the cluster of the origin.}
Fix $p > \pc(0)$, so that $\mathbb P_{p,0}(|\mathcal C_0| = \infty) > 0$. Throughout this proof, we abbreviate~$\alpha = \alpha(p,0)$ and~$\beta = \beta(p,0)$.

Note that for any $n\in\bbN$
\begin{align*}
\max \xi_n(\{0\}) &{}\le r_n,&\min \xi_n(\{0\}) &{}\ge l_n
\end{align*}
and on the event $|\cC_0|=\infty$, for all $n\in\bbN$,
\begin{equation}
\label{eq_note_cluster}
l_n = \min \xi_n(\{0\}) \le \max \xi_n(\{0\}) = r_n,
\end{equation}
{using the non-crossing property of simple paths.} Taken together with~$\frac{l_{2n}}{2n} \xrightarrow{n \to \infty} \beta$ and~$\frac{r_{2n}}{2n} \xrightarrow{n \to \infty} \alpha$,~\eqref{eq_note_cluster} implies that~$\beta \le \alpha$.

For~$a > 0$, we write
\[V_-(a):= \{(v,n) \in \mathbb Z \times 2 \mathbb N:\; v \le an\},\quad V_+(a):=\{(v,n) \in \mathbb Z \times 2 \mathbb N:\; v \ge an\}.\]
 We claim that
\begin{equation}
\label{eq_note_alpha}
    \alpha = \inf\{a > 0:\; \mathbb P_{p,0}(|\cC_0 \cap V_+(a)| < \infty) = 1\}.
\end{equation}
To see this, first take~$a > \alpha$. Since~$\max \xi_n(\{0\}) \le r_n \; \forall n$ and~$\frac{r_{2n}}{2n} \xrightarrow{n \to \infty} \alpha$, we see that almost surely there are only finitely many~$n \in 2 \mathbb N$ such that~$\max \xi_n(\{0\}) \ge an$, so there are almost surely only finitely many points in~$\cC_0 \cap V_+(a)$. On the other hand, if~$a < \alpha$, we have
\begin{align*}
\mathbb P_{p,0}(|\cC_0 \cap V_+(a)| = \infty) &\ge \mathbb P_{p,0}(\max \xi_n(\{0\}) \ge an \text{ for infinitely many } n \in 2 \mathbb N) \\
    &\ge \mathbb P_{p,0}(|\cC_0| = \infty) > 0,
\end{align*}
where the second inequality follows from~\eqref{eq_note_cluster} and~$\frac{r_{2n}}{2n} \xrightarrow{n \to \infty} \alpha$. This concludes the proof of~\eqref{eq_note_alpha}. Similarly, we have
\begin{equation}
    \label{eq_same_beta}
\beta = \sup\{b > 0:\; \mathbb P_{p,0}(|\cC_0 \cap V_-(b)| < \infty) = 1\}.
\end{equation}
Let~$\Phi: \mathbb R^2 \to \mathbb R^2$ be the reflection about the diagonal~$y=x$, that is,~$\Phi(x,y) = (y,x)$. Since we are taking~$q = 0$, our model has the symmetry~$\cC_0 \stackrel{\text{(law)}}{=} \Phi(\cC_0)$. In particular, for any~$a > 0$,
	\begin{align*}
		\mathbb P_{p,0}( | \cC_0 \cap V_-(1/a) | < \infty) =\mathbb P_{p,0}( |\cC_0 \cap V_+(a)| < \infty).
	\end{align*}
	Together with~\eqref{eq_note_alpha} and~\eqref{eq_same_beta}, this gives~$\beta = 1/\alpha$. We already had~$\beta \le \alpha$, so we obtain~$\beta \le 1$ and~$\alpha \ge 1$.
\end{proof}

\end{document}
